% Page 043

\seite{43}

\abschnitt{Spaziergang [?]}
\pstart
Am 25.\ August machten die mit den Schülern der\ob
Pfarrei Seffern\unsicher{} gemeinschaftlich einen\ob
Spaziergang nach Kyllburg\unsicher.

\pend
\abschnitt{Ernte}
\pstart

Die diesjährige Obst- und Getreideernte lieferte\ob
einen reichlichen\unsicher{} Ertrag. Bei günstiger\ob
Witterung konnten die Früchte gut eingebracht\ob
werden. Ein großer Theil der Futterkräuter\unsicher\ob
wurde wegen Mangel an Raum auf Haufen\ob
aufgestellt. Die Witterung war eine recht günstige,\ob
weshalb alles rechtzeitig und trocken eingefahren\ob
werden konnte. Der Haferertrag ergab nur\ob
einen mittleren Ertrag. Obst gab es nur wenig. Heu die\ob
gleichen\unsicher{} Erträge\unsicher.

\pend
\abschnitt{Kirchenfeierlichkeiten}
\pstart

Der Geburtstag Sr. Maj. des Kaisers wurde am\ob
27.\ Januar\unsicher{} zu Seffern gemeinschaftlich von den\ob
vier Orten der Pfarrei unter Leitung der Orts\ohb
schulgehülfen\unsicher{} gefeiert. Abwechselnd mit Liedern,\ob
die von den größeren Kindern gesungen, folgten\ob
Gedichte, die von einzelnen Kindern vorgetragen\ob
wurden. Die Ansprache an die Kinder hielt\ob
Herr Lehrer Klein\unsicher{} aus Bickendorf.

\pend
\abschnitt{Revision}
\pstart

Die hiesige Schule wurde am 25.\ Februar durch den\ob
Herrn Kreisschulinspektor Lenz\unsicher{} aus Bitburg\ob
revidiert.

\pend
\abschnitt{Osterprüfung}
\pstart

Die Osterprüfung wurde am 18.\ März durch den\ob
Herrn Pfarrer Schmitt aus Seffern abgehalten.

\pend
\abschnitt{gf. 18/3.     Schmitt, Pfr}
\pstart


\pend
\jahresueberschrift{Das Schuljahr 1907-08}
\pstart

Entlassen wurden in diesem Jahre 5 Knaben und\ob
3\unsicher{} Mädchen. Die meisten widmen sich der\ob
Landwirthschaft. Aufgenommen\unsicher{} Knaben.\ob
Die Schülerzahl beträgt, da auch ein fremdes\unsicher\ob
Kind die hiesige Schule besucht,\ob
15 Knaben und 11\unsicher{} Mädchen. Ein Knabe ist\ob
nicht ganz willfährig.
\pend
